\chapter{Hearing Loss}

\section*{Definition}
A partial or total inability to hear sound in one or both ears, ranging from mild (difficulty with soft speech) to profound (deafness), classified as conductive (outer/middle ear), sensorineural (inner ear/auditory nerve), or mixed type.

\section*{What Happens in Your Body}
Conductive hearing loss results from obstruction or damage to the external ear canal, tympanic membrane, or ossicular chain impeding sound transmission to the cochlea. Sensorineural hearing loss involves damage to cochlear hair cells, the auditory nerve, or central auditory pathways. Noise exposure causes mechanical damage to hair cells and metabolic exhaustion. Presbycusis (age-related) involves cumulative hair cell loss, stria vascularis atrophy, and central processing decline.

\section*{Causes}
\begin{itemize}
  \item Presbycusis (age-related hearing loss, most common)
  \item Noise-induced hearing loss (occupational, recreational, acute acoustic trauma)
  \item Otitis media with effusion (glue ear, most common in children)
  \item Otosclerosis (abnormal bone remodeling of stapes)
  \item Cerumen impaction (earwax blockage)
  \item Meniere disease (episodic vertigo, tinnitus, fluctuating hearing)
  \item Ototoxic medications (aminoglycosides, loop diuretics, cisplatin, salicylates)
  \item Sudden sensorineural hearing loss (viral, vascular, autoimmune)
  \item Acoustic neuroma (vestibular schwannoma)
  \item Trauma (temporal bone fracture, barotrauma, perforated TM)
\end{itemize}

\section*{When to See a Doctor}
See your doctor if:
\begin{itemize}
  \item Symptoms are severe or getting worse over time
  \item Sudden hearing loss (within 72 hours), hearing loss with vertigo, tinnitus, or ear pain, or hearing loss in one ear only
  \item You have other concerning symptoms such as fever, weight loss, or bleeding
\end{itemize}

\section*{Related Symptoms}
\begin{itemize}
  \item Tinnitus
  \item Earache
  \item Vertigo
  \item Ear fullness
  \item Balance problems
\end{itemize}
\textbf{Important:} If this symptom persists, your doctor will check for serious underlying causes.

\section*{Self-Care / Home Management}
\begin{itemize}
  \item Protect your ears from loud noise by wearing earplugs or noise-canceling headphones.
  \item Keep the volume at 60\% or less when using headphones or earbuds.
  \item Have earwax removed by a healthcare professional — do not use cotton swabs.
  \item Avoid ototoxic medications (aspirin, NSAIDs in high doses) unless prescribed.
  \item Seek prompt medical evaluation for sudden hearing loss, as early treatment improves outcomes.
\end{itemize}

\section*{How Your Doctor Will Evaluate You}
\begin{itemize}
  \item Your doctor will examine the ear canal and tympanic membrane with an otoscope.
  \item Tuning fork tests (Rinne and Weber) help differentiate conductive from sensorineural hearing loss.
  \item Pure-tone audiometry and speech audiometry quantify the degree and type of hearing loss.
  \item Tympanometry assesses middle ear function and eustachian tube status.
  \item MRI of the internal auditory canals may be ordered if acoustic neuroma is suspected.
\end{itemize}

\section*{Treatment Options}
\begin{itemize}
  \item Hearing aids are the mainstay of treatment for age-related and noise-induced hearing loss.
  \item Removal of cerumen impaction for conductive hearing loss from earwax.
  \item Oral corticosteroids for sudden sensorineural hearing loss, ideally within 72 hours of onset.
  \item Surgical intervention (stapedectomy, tympanoplasty, cochlear implantation) for appropriate candidates.
  \item Avoiding ototoxic medications and using hearing protection to prevent further damage.
\end{itemize}

\section*{References}
\begin{thebibliography}{99}
\bibitem{presbycusis1} Gates GA, Mills JH. "Presbycusis." Lancet. 2005;366(9491):1111-1120.
\bibitem{presbycusis2} Yamasoba T, Lin FR, Someya S, et al. "Current concepts in age-related hearing loss: epidemiology and mechanistic pathways." Hear Res. 2013;303:30-38.
\bibitem{presbycusis3} Cunningham LL, Tucci DL. "Hearing loss in adults." N Engl J Med. 2017;377(25):2465-2473.
\bibitem{presbycusis4} Lin FR, Metter EJ, O\'Brien RJ, et al. "Hearing loss and incident dementia." Arch Neurol. 2011;68(2):214-220.
\bibitem{presbycusis5} Sprinzl GM, Riechelmann H. "Current trends in treating hearing loss in elderly people." Gerontology. 2010;56(3):351-358.
\bibitem{presbycusis6} Frisina RD, Walton JP. "Age-related hearing loss." In: Cummings Otolaryngology. 7th ed. Elsevier; 2021.
\end{thebibliography}
